% Esqueleto didático adaptado do pré-projeto fornecido pelo aluno.
\section{Planejamento da análise e discussão}

A etapa de preparação dos dados será essencial para garantir a consistência da base analítica. Inicialmente, serão realizadas atividades de extração, padronização, limpeza, tratamento de inconsistências e integração das diferentes fontes. No caso dos dados institucionais, será necessário padronizar nomes de cidades, aeroportos, companhias aéreas, datas, valores, status de passagens e demais campos relevantes. Também será necessário identificar registros duplicados, valores ausentes, passagens canceladas, remarcadas ou reembolsadas, bem como eventuais divergências entre data de solicitação, autorização, emissão e viagem.

Também poderão ser exploradas variáveis institucionais, como área demandante, finalidade da viagem, prazo entre solicitação e emissão, urgência, histórico de remarcação e padrões recorrentes de deslocamento, desde que disponíveis e adequadas ao uso. Essas variáveis poderão contribuir para adaptar o modelo às restrições reais de compra em ambientes institucionais.
